% Reviewable insertion, not a patch and not included automatically by article.tex.
% Base commit: 4c173c06cc32c9cea554b39be1837ad2ae897fc1.
% Include only after the corresponding proof sections and bibliography entries
% have been integrated, so the ker:/ord:/mesh:/neg:/mix: references resolve.
% This text changes status and scope; it does not replace the proofs.

\subsection*{Status updates from the extremizers and zero-threshold continuation}

The following updates refer to the immutable ProveIt snapshot
\texttt{4c173c06cc32c9cea554b39be1837ad2ae897fc1} and its six incoming
archives. They record ordinary proved results and precise remaining
questions; they are not a claim of proof-assistant formalization or an
exhaustive priority claim over the literature.

\paragraph{Pointwise Euler-kernel maxima.}
The global limiting-kernel conjecture labelled \texttt{conj:limit}
in \cite[Section 8]{IncomingSharp} is proved by
Theorem~\ref{ker:thm:limit}, including the unique limiting maximizer
and the global control missing from the earlier compact scaling limit.
The all-finite-index conjecture \texttt{conj:maxima} is only partly
settled: Theorem~\ref{ker:thm:N2} proves the case $N=2$, while
Theorem~\ref{ker:thm:analytic} proves eventual uniqueness, strict decrease,
and a convergent expansion. The assertion at every intermediate finite
index remains open.

\paragraph{Optimal averaged Euler constants.}
The conjecture \texttt{conj:euler-rate} in
\cite[Section 9]{IncomingBounds} is proved. With
\[
 p=\frac{\log2}{\log(3/2)},\qquad
 q=\frac{\log3}{\log2},\qquad c=(1-q^{-1})q^{-p},
\]
Theorem~\ref{ord:thm:leading} gives $C_N-1\sim cN^{-p}$.
The global step is Theorem~\ref{ord:thm:axis-reduction}: for all
sufficiently large $N$, the positive-order supremum equals the attained
maximum on the outer-order axis. Positive outer orders do not attain
that value. This exact reduction also proves eventual strict decrease
of $C_N$. The axis maximizer is eventually unique.

Theorem~\ref{ord:thm:second} adds the threshold correction
$DN^{-\kappa}(\log N)^{-1/2}$ and its inverse-log expansion.
Proposition~\ref{ord:prop:order-profile} and
Corollary~\ref{ord:cor:near-maximizers} give the exact criterion for an
$o(C_N-1)$ near maximizer. These statements should retain their accuracy
scale and eventual quantifiers. No moderate-index error guarantee or
numerical starting index is supplied by the asymptotic theorem.
The pointwise constants $M_N$, averaged constants $C_N$, and the already
proved all-index universal constant $C_*$ remain distinct.

\paragraph{Polynomial reciprocal-gamma Appell separation.}
The polynomial-separation question in
\cite[Sections 4 and 8]{IncomingUniform} is answered by
Theorem~\ref{mesh:thm:polynomial}:
\[
 \operatorname{mesh}(Q_n)\ge\frac1{10n^2}\qquad(n\ge2).
\]
Retain the earlier exponential mesh bound; it is compatible with the
new result and may be stronger at a small degree. The effective zero
framework is strengthened by Theorems~\ref{mesh:thm:saturation} and
\ref{mesh:thm:transfer}. In particular,
\[
 K_n=\left\lceil(16+240n^2H_{n-1})^2\right\rceil
\]
is a sufficient derivative threshold, uniformly for $0\le\rho\le1$,
and the simultaneous logarithmic location error is less than
$8/\sqrt k$. This is not an optimal-threshold assertion. At $\rho=1$
the theorem concerns the derivative family $k\ge1$, not the generally
divergent undifferentiated Lerch series.

\paragraph{The first two negative outer orders below minus one.}
Keep the prior negativity theorem for every real outer order at least
$-1$, and retain the provenance of the previously evaluated points.
Theorem~\ref{neg:thm:transitions} proves that for each $m=2,3$ the
function $b\mapsto g_{-m,b}$ has exactly one positive zero, which is
simple, with positive sign before it and negative sign after it.
The two zeros satisfy $0<b_2<1<b_3$. No all-negative-order one-crossing
claim follows from this result.

\paragraph{Exact mixed-generator identities and unresolved periods.}
Theorems~\ref{mix:thm:shift}, \ref{mix:thm:rational}, and
\ref{mix:thm:integer}, together with
Corollary~\ref{mix:cor:resummation}, add exact shift, rational sampling,
parameter-jet, integer-sampling and finite-part formulas for the mixed
Gaussian generator. The resulting convergent identities combine weights;
they do not prove the individual unresolved fixed-weight reductions.

The $S_4$ formula is already a theorem. Its original aliases
\texttt{s4proof:thm:s4}, \texttt{gaussian:thm:S4}, and
\texttt{gaussian:conj:S4} should retain that status, notwithstanding the
historical name of the last alias. By contrast, the $S_6$ and $S_8$
statements labelled \texttt{cycloquot:conj:S6} and
\texttt{s8new:conj:S8}, and the incoming $S_{10},S_{12}$ vectors,
remain conjectural. Conjecture~\ref{mix:conj:S14} appends a frozen
candidate of total weight fifteen. Proposition~\ref{mix:prop:proximity}
proves the normalized residual bound $|R_{14}|<10^{-775}$ by exact
rational arithmetic. The delivered interval contains zero; exact
vanishing is still conjectured.

\paragraph{Evidence and integration scope.}
The source identities and hashes are recorded in the companion
provenance. No theorem-level error was established in the inherited
inputs used here; the changes above are status updates and scope
qualifications. Restricted relation-space obstructions are not claims
of period independence. The analytic proofs, finite exact certificates,
and numerical diagnostics should continue to be distinguished when
these sections are incorporated into the manuscript.
